\chapterNote{Based on a talk given on November 27, 2020}

\chapter{Turning Toward a True Refuge}
\tocChapterNote{Based on a talk given on November 27, 2020}
%\tocChapterNote{November 27, 2020}

\looseness=-1
This morning, we received a teaching on the contemplation of the five khandhas, and in particular, the \emph{upādānakkhandhas}, the aggregates affected by \emph{upādāna} (clinging). The reflection included the kinds of clinging, different descriptions of the khandhas, similes that the Buddha used, and the importance of contemplating them through the lenses of \emph{anicca} (impermanence), \emph{dukkha} (suffering), and \emph{anatta} (not-self).

Contemplation of the aggregates should be one of our regular meditation practices as it's at the crux of the path to freedom. It deepens our understanding of what we usually regard as me or mine. Developing insight into the process of I-making and mine-making will lead to seeing the bald-faced reality of what it is that we have been taking refuge in. This understanding leads to a sense of disenchantment and supports relinquishment. This is the path to freedom. Freedom is not found in becoming anything special or having all sorts of high consciousness experiences. It comes from seeing through our presumptions of stability, close examination of where we seek pleasure, and most importantly, contemplation of the perception of self. We look closely at what we are taking refuge in, questioning the validity of that refuge as a long-term source of happiness. We need to see the truth of our reality to the point where we want to turn away from that false refuge and see something deeper.

\looseness=-1
Most of us have been doing this routine of pleasure-seeking long enough to have a sense that it's not taking us exactly where we want to go. This realization is the impetus, the inspiration for seeking something beyond this experience of me trying to be happy, me trying to find satisfaction in the world. We are looking to actualize the transcendent goal of freedom instead of constantly striving for \mbox{temporary happiness.}

In the suttas, there is a discussion between Venerable Sāriputta and Venerable Mahākoṭṭhita that emphasizes reflection on the five khandhas \mbox{(SN 22.122)}. The discussion of what to look for in examining the five khandhas is found in many different suttas and includes taking stock of the five khandhas, form, feeling, perceptions, mental formations, and consciousness, and experiencing them through the lenses of anicca, dukkha, and anatta.

Several other words are used to describe how to regard the five khandhas, and they have a powerful impact when we hear them recited. We view each of the five khandhas in turn, starting with form, as impermanent, suffering, a dart, a tumor, an affliction, an illness, a misery, a disintegration, alien, empty, and not-self. These are very strong words. Similarly, when we experience and examine the khandhas of feeling, perception, mental formations, or consciousness, we see them all in the light of impermanence, suffering, a tumor, a dart, an illness, an affliction, a misery, disintegrating, empty, and not-self. These words grab our attention and give rise to alert reflection.

As a result of reflecting on these words, we can ask, ``What am I doing taking refuge in the five khandhas?'' Disenchantment arises from the realization of where we have put all of our hopes and attention.

We can start to reflect deeply on the five khandhas, to examine them and to see if we can uproot the perception of self in relation to them. The perception of self comes in different forms: as ``I am,'' ``I am this body,'' or ``I am these feelings, perceptions, thoughts, or consciousness.'' The notion of self can also manifest as a perception of ownership, the experience that they are mine. ``This is my body. There is somebody here who owns this body, who possesses these feelings, perceptions, thoughts, or consciousness.'' We may not view them as ``me'' or ``I am some part of these or all of these,'' but there still can be a subtle view of somebody who is here, who possesses them, owns them and should have some control over them. That's what a self would do. That's what a ``me'' would do if it were solid, real, and in control.

We need to pick apart self-view, experience it, and question it closely, which is what our settling and calming meditation is designed to support. We aim to get clear, steady, contented, and peaceful enough to be able to hold a contemplation like this, to be with it for a while and delve into it. It doesn't mean we can't think about it in a skillful way before our minds collect and pull together. In fact, engaging in active contemplations can help settle the mind even more. We talk about \emph{samatha} (calm) and \emph{vipassanā} (insight) being two ends of the same stick. They help develop and support each other.

Contemplation of the body encourages exploring the body in relation to anicca, dukkha, and anatta and leads to questioning the assumptions about the nature of the body and our firmly entrenched ways of viewing ourselves. There are many different ways that we can contemplate the body, and it helps to start with a firm determination to make this our Dhamma lens, our focus for exploration.

Earlier this year (2020), when I stepped down from being co-abbot with Ajahn Nyāṇiko, he put together a lovely book of photographs from throughout my life and personal recollections from many people, mostly from my monastic life. It's a great collection. One of the photographs he got from my sister was of me when I was less than two years old, a photo I had never seen before. My mother was holding me in her right arm, balanced off to one side, and my brother, sister, grandmother, and a young neighbor friend were in the photo. I had a slight scowl on my face and looked like I wanted to get down. When I look at that photo, the question comes up, ``Was that me?'' It doesn't seem like me now. It was a touching image, but it brings up the question, ``Why would I call that me, when it feels like this me is so different from that me?''

There was another photo of the Sanghapāla Board from the time before I was ordained. It was before Abhayagiri Monastery was formed and was during one of Ajahn Amaro's trips to the US, when he stayed for a few months. Towards the end of this visit, there was a Sanghapāla Board meeting, and although I wasn't on the Sanghapāla Board, I came down from Seattle to join the meeting. I was in my late 30s, as were Debbie and Ajahn Amaro, who were both in the photograph. The three of us were sitting there looking quite young and thin. As I was looking at that photograph, I thought, ``Well, Ajahn Amaro looks pretty much the same. He has aged some and put on a little bit more weight, but Debbie and I look substantially different.'' In the photo, I had brown hair, a beard, and was wearing trousers and a sweater. When I look in the mirror now, I see almost no hair, certainly no beard, and what hair is present is white instead of brown. I'm clearly not in my late 30s anymore.

\looseness=-1
The photos were of when I was two years old and in my late 30s, but I am now in my mid 60s. Is this the same person? Biologically, we know that human body cells change out more or less every seven years, some much faster. Some turn over on an almost daily basis, and others are quite long lasting. Anatomically, we're not the same being that we were years ago. Where did that ``me'' go? Most of the cells that were me, that made up this body seven years ago, have died and have been excreted through urine or feces. They have been replaced with similar ones, and the body in general has morphed and aged. There is merely an illusion of continuity that carries through.

The same is true for all feelings, perceptions, and memories. Where did Tom DeMaria go? Where did Sāmaṇera Karuṇadhammo go? These are all just creations in our minds, even the Karuṇadhammo Bhikkhu of yesterday, who was sitting in this spot answering questions. Where did he go? It's just a memory. Everything prior to this moment is just a memory. The same is true for everyone. We can reflect on our assumptions and views of who we are and realize that this perception of self is a manufactured illusion to keep a sense of continuity in our lives. But it's a construct. It's an assumed reality. Realizing this truth can be quite a relief.

It's easy and so tempting to hang on to memories and make them into present realities. Many memories are painful, and hanging on to them leads to more and more dukkha. Instead, it is a relief to be able to say, ``Oh, that event is finished and is only a memory. I don't need to impute any particular reality on to it other than as a memory.'' And then here I am now. Here, just as it is. The future is a guess, a projection. There really isn't anything happening other than what's happening right here and now.

Reflecting on the illusion of continuity is one way we can start to get a handle on the perception or assumption of self. Another way to view it is by looking through the lens of control. We can investigate what we have control over and what we don't have control over, and how \mbox{that affects self-view.}

There is a modicum of control relative to our bodies and minds. Intentions arise to perform an action through either body, speech, or mind. An intention arises based on some cause or condition, such as sensory contact or a memory of an event. Following the moment of contact, an underlying tendency of desire or aversion gives rise to an intention to do something, to act or speak.

But is there really an independent agent, a self, who is in control? The experience of intention arising is often filled with the notion of me or mine, which imbues the action with a perception of control. We do have the ability to control our experience to a limited extent. For example, I can form the intention to raise my right hand, and then my hand goes up. That means I'm in control of my hand, right? I do what I want to do with it when I want to do it, but it only goes so far. If we reflect on the body, we notice it follows a natural course of experiences which we don't have any control over. The aging process keeps on going. We can say to ourselves, ``Okay, body, don't age!'' but it doesn't pay any attention to us. The intention falls flat. There is no control in this case. Old age, sickness, and death are inevitable. We can influence our health, modify it with the food we eat, the medicines we take, or the exercises we engage in. But in reality, we're all in different places on the same conveyor belt, and it's moving in one direction. We need to closely investigate the sense of control. Who's in charge of this body? Nature is in charge. In quiet moments, we can have this insight when we deeply investigate the notion of control relative to the body.

\looseness=-1
Right now, sitting quietly, compose the mind and tune into the presence of the body. Allow the body to relax, and keep your focus on the whole body, not letting your attention wander around. Start to notice the various sensations, feelings, and sensitivity of the body. A phrase that one of our teachers uses for this bodily presence is ``intelligent sensitivity.'' This body is one of intelligent sensitivity. It has its own form of intelligence and knows what it needs to do. It will mostly keep on doing what \mbox{it needs to do in} the best way it can, oftentimes in opposition of what we intend. You can notice that when you are sitting still. Notice what are called autonomic functions of the body. The autonomic nervous system is that which controls the breathing, heartbeat, digestion, and other bodily processes. These keep on functioning, regardless of our input and wishes. You think you should have control, but tune into the breathing process for a moment and ask yourself how much control you have. You can command yourself to hold your breath for a short while, but how long does that last? It can be interesting to hold your breath for as long as you can with the intention of having control. After ten or twenty seconds, or a little bit longer, the body overrides any intention that you have and will gasp and take a breath. It's just what bodies do.

Alternatively, tune into your heartbeat. You can feel the heart rate speed up a little bit through intense mental activity or when you are excited or angry. You might be able to slow it down consciously during meditation. However, the heart keeps beating sixty to eighty times a minute, every hour, twenty-four hours a day, seven days a week, 365 days a year, for however many years you live. The heart has an intelligence that is not a cognitive intelligence, but a bodily intelligence. It knows what to do whether you're paying attention to it or not.

You can also pay attention to other sensations, such as slight muscle movements. You're probably shifting position frequently in some way or another, maybe in a very small way. You're not consciously intending every movement, the body is just doing it. Be with the processes of subtle, twitching muscles of the face, or of swallowing saliva. Scan through the body and sense the presence of the body and its movements, its energies, its elements. Gently ask yourself, ``Who's in charge?'' In actuality, no one is. What a relief!

\looseness=-1
You can reflect on the nature of control in regard to the other aggregates. For example, you can contemplate the pleasant and unpleasant feelings that arise and fade away. The Buddha compared feelings to bubbles forming and bursting in a water puddle when it rains. Feelings arise and pass away, one after another, after another, after another. Sense impressions, memories, and perceptions all do the same thing. Likewise with thoughts and emotions: Where did that mood go that you were wrapped up in yesterday or the day before? Those strings of stories and thoughts, where are they now? Do you really have control or are the defilements producing them? Greed, anger, irritation, and the different forms of confusion or anxiety are generating thoughts and feelings. There's not anybody in there doing it.

\looseness=-1
There is also the good stuff. By putting attention on developing generosity, kindness, patience, and other skillful mental qualities, pleasant wholesome thoughts, feelings, and emotions will result. But, again, it's not ``me'' \mbox{who's doing it.}

Lastly, the consciousness aggregate weaves together its story of our world by jumping from one sense base to another and creating a story, but there is no independent, self-existent being doing it. This aspect of experience is only a series of perceptions and feelings, which lead to intentions, thoughts, or bodily actions, both skillful and unskillful, which circle around to produce more perceptions and feelings.

We can use the quiet, contemplative mind to question who we are, what we assume about our reality, and what we assume about ourselves. It is all just an assumed and constructed reality that seems very real when we're engrossed in it, involved with it, and acting on it. But when we take time to look and see these mental activities just for what they are, the fascination with what we can get from them lessens. We can gratify a desire or a craving, such as a craving for being or existence, but these are just fleeting desires, finding expression through the five khandhas. We can use the power of a calmed and settled mind for strategic questioning of who we think we are.

\looseness=-1
Similar to the discussion regarding how to view the five aggregates mentioned earlier between Sāriputta and Mahākoṭṭhita, there is a sutta in the Aṅguttara Nikāya in which the Buddha instructs a disciple in a similar way (AN 9.36). While engaged in the first \emph{jhāna}, the disciple is to regard any phenomena that arise in relation to the five khandhas as impermanent, suffering, a dart, a tumor, an affliction, a misery, disintegration, alien, empty, not-self. The Buddha then goes further and instructs the disciple to turn away from those phenomena and instead direct the mind to the deathless element. The instruction to ``direct one's mind to the deathless element'' occurs in several places in the teachings and does not always include an explanation (for example, AN 3.131). How does one do that? If one hasn't experienced the deathless element directly, how does one know how to direct one's attention to it? But this particular sutta in the Aṅguttara Nikāya continues with, ``Direct one's mind to the deathless element thus: `This is peaceful, this is sublime, that is, the stilling of all formations, the relinquishment of all acquisitions, the destruction of craving, dispassion, cessation, Nibbāna.'\,'' This is a description of how to directly contemplate the deathless element. But if the mind is not yet ready for the direct experience, this instruction can be used as a perceptual training to imagine what the peace of Nibbāna would be like. Simply reciting these phrases can be a settling and insightful contemplation.

\looseness=-1
With the dispassion that results from seeing the khandhas as impermanent, unsatisfactory, and not-self, we are no longer enthralled with the temporary gratification that the khandhas offer. We can incline towards reflecting on a much deeper peace, the peace of Nibbāna. This is peaceful, this is sublime. The stilling of all formations. That's when the kammic process comes to a standstill. We are not creating any new kamma. Formations (saṅkhāras), based on mental volition, take a pause. Volitions based on greed, hatred, and delusion stop. ``The relinquishment of all acquisitions (\emph{upadhi}).'' ``Acquisitions'' sometimes refers to the five aggregates affected by clinging, which are entwined with the perception of self as ``me'' or ``mine.'' We cling to the five khandhas with a sense of ownership or identification. Sometimes upadhi are also called the bases or constituents of existence. With disenchantment and dispassion, we relinquish them. There is no need to hold on to them any longer, no need to keep that construct going, no need to cling anymore.

Deep peace also comes with the destruction of craving, ``\emph{taṇhā}''. Craving is the momentum, the mental push to get something from the sensory world. It is also the momentum to be somebody, to become something, to have an identity, or simply to exist. It is also the craving to annihilate, the aversion to dealing with anything unpleasant. With the destruction of craving, dispassion is experienced. We lose the passion for anything related to the five khandhas, anything related to the conditioned realm. It doesn't mean we cease to exist, but the passion for that which no longer fulfills just fades away. \emph{Nirodha} is translated as cessation or non-arising. Sāriputta said, ``Nibbāna is the cessation of becoming.'' It is the non-arising of the urge to become, to exist. That kind of non-arising is still, peaceful. Dispassion, cessation, Nibbāna: cooling, unbinding, complete freedom, the end of suffering.

\looseness=-1
This is where contemplation can take us. If we have yet to arrive at that point of deep realization, then bringing these phrases to mind as a possibility, as a perception, is a very powerful contemplation. The teachings of the Buddha and the words that he used to describe the path and its results are powerful, profound, and provide inspiration to \mbox{keep on the path.}
